% ATTEMPT_STATUS: unresolved
% ATTEMPT_ORIGIN: new research, 2026-10-01
% TOP_PROBLEM: {"id": 30, "title": "The Lonely Runner Conjecture", "category_id": 2, "category": {"id": 2, "name": "combinatorics", "display_name": "Combinatorics", "description": "Counting problems, graph theory, discrete structures.", "slug": "combinatorics", "order_index": 2, "created_at": "2026-07-31T15:26:25.671Z"}, "status": "open"}
% FIRST_POSTED: 2026-10-01
% Imported from: unsolvedmath_additional_500.tex; UnsolvedMath ID 30
% !TeX program = lualatex
% Compile twice: lualatex 30.tex
% Self-contained: no external bibliography, images, or \input files.
% Numeric section labels and visible numbers are original UnsolvedMath IDs.
\documentclass[11pt,a4paper]{article}
\usepackage[margin=24mm,headheight=14pt]{geometry}
\usepackage{fontspec}
\setmainfont{Latin Modern Roman}
\setsansfont{Latin Modern Sans}
\setmonofont{Latin Modern Mono}
\usepackage{amsmath,amssymb,mathtools}
\usepackage{unicode-math}
\setmathfont{Latin Modern Math}
\usepackage{microtype}
\usepackage{enumitem,xurl,needspace}
\usepackage[dvipsnames]{xcolor}
\usepackage{hyperref}
\hypersetup{unicode=true,colorlinks=true,linkcolor=MidnightBlue,urlcolor=MidnightBlue,
 pdftitle={UnsolvedMath Problem 30},pdfauthor={Source-annotated compilation},bookmarksnumbered=true}
\usepackage{bookmark,tocloft,titlesec,fancyhdr}
\setlength{\cftsecnumwidth}{4.2em}
\cftsetpnumwidth{2.5em}
\cftsetrmarg{4em}
\titleformat{\section}{\Large\bfseries\raggedright}{\thesection}{0.7em}{}
\pagestyle{fancy}\fancyhf{}
\fancyhead[L]{\small\sffamily UnsolvedMath research attempt}
\fancyhead[R]{\small Original catalogue IDs}
\fancyfoot[C]{\thepage}
\setlength{\parindent}{0pt}
\setlength{\parskip}{5pt plus 1pt}
\setlength{\emergencystretch}{3em}
\setcounter{tocdepth}{1}\setcounter{secnumdepth}{1}
\setlist[itemize]{leftmargin=3em,itemsep=4pt,topsep=4pt,parsep=0pt}
\allowdisplaybreaks
\newcommand{\N}{\mathbb{N}}
\newcommand{\Z}{\mathbb{Z}}
\newcommand{\Q}{\mathbb{Q}}
\newcommand{\R}{\mathbb{R}}
\newcommand{\C}{\mathbb{C}}
\newcommand{\F}{\mathbb{F}}
\newcommand{\A}{\mathbb{A}}
\newcommand{\PP}{\mathbb{P}}
\newcommand{\E}{\mathbb{E}}
\newcommand{\eps}{\varepsilon}
\newcommand{\dd}{\,\mathrm d}
\newcommand{\sref}[2]{\hyperlink{source:#1:#2}{[S#2]}}
\newif\ifshowcatalogue
\showcataloguetrue
\newcommand{\cataloguescope}[1]{\ifshowcatalogue
 \paragraph{Statement recorded in the supplied source.}
 #1\fi}
\begin{document}
% UnsolvedMath ID 30; source catalogue code COMB-004

\setcounter{section}{29}

\section{The Lonely Runner Conjecture}\label{30}

{\small\sffamily UnsolvedMath ID 30 · Combinatorics}

\subsection{Definitions and mathematical statement}

\paragraph{Shared notation.}Unless a section says otherwise, $\N=\{1,2,\ldots\}$,
$\Z$ is the ring of integers, $\Q,\R,\C$ are the rational, real, and complex
fields, $[N]=\{1,\ldots,\lfloor N\rfloor\}$, and $\log$ is the natural
logarithm. For real functions of a parameter tending to infinity,
$f\ll g$ and $f=O(g)$ mean $|f|\leq Cg$ eventually for some fixed $C$;
$f\gg g$ reverses this comparison; $f=o(g)$ means $f/g\to0$;
$f\sim g$ means $f/g\to1$; and $f\asymp g$ means both order bounds.
A subscript on $O$ or $\ll$ specifies allowed constant dependence.
The expression $n^{o(1)}$ in an upper bound means growth smaller than
$n^\epsilon$ for each fixed $\epsilon>0$. Local definitions govern any
exception, finite-field convention, or use of the same symbol for another
object. An asymptotic claim is not automatically an effective algorithm.



Use the unit circle $\R/\Z$, distance $\|u\|_{\R/\Z}=\min_{m\in\Z}|u-m|$, and the standard common-start model $x_i(t)=v_it\pmod1$. The speeds $v_1,\ldots,v_n$ are distinct real numbers. For each $i$, the conjecture asks for a time $t>0$ such that $\|(v_i-v_j)t\|_{\R/\Z}\geq1/n$ for every $j\ne i$. A different time may serve each runner. \sref{30}{1} \sref{30}{2}

\paragraph{Statement recorded in the supplied source.}
For any $n$ runners on a circular track with distinct constant speeds, each runner is "lonely" (distance at least $1/n$ from all others) at some time. \sref{30}{1}

\paragraph{Residual task reported by the archive.}

The embedded review (2026-08-17) records: Prove the conjecture for every number of runners. \sref{30}{1}

\subsection{Short English statement}

Starting together on a circular track, must each runner eventually be separated from every other runner by at least one nth of the circumference?

\subsection{Research attempt}
\textbf{Status: UNRESOLVED.} We prove a universal weaker separation bound for arbitrary real relative speeds, establish the exact conjectured bound for certain congruence classes of rational speed systems, and give a finite exact optimization reduction for any fixed integer tuple. These results do not prove the required bound uniformly over all tuples.

\paragraph{Literature and scope check.}
The primary abstract of Sungkawichai--Trakulthongchai [S2], version 2 consulted 1 October 2026, reports a computer-assisted proof for ten, eleven, and twelve nonzero integer speeds, corresponding to eleven through thirteen runners. Those computations were not reproduced here. The target in this file quantifies over all numbers of runners and real speeds. We therefore keep our elementary argument separate from the cited finite-runner results.

Fix the runner being isolated. Subtract its speed from all speeds and take absolute values of the others: $u_j=|v_j-v_i|>0$, $1\le j\le k=n-1$. Sign changes are harmless because $\|-x\|=\|x\|$. Equal absolute relative speeds can occur, although the original signed speeds are distinct; none of the arguments below assumes the $u_j$ are distinct. Define $F(t)=\min_j\|u_jt\|$. The target is $F(t)\ge1/(k+1)$ at some positive time.

\paragraph{A real-speed bound from measure.}
For $0<\delta<1/2$ put $B_j=\{t\ge0:\|u_jt\|<\delta\}$. This set is periodic with period $1/u_j$ and occupies fraction $2\delta$ of each full period. Cutting $[0,T]$ into full periods and a last remainder gives
\[
|B_j\cap[0,T]|\le 2\delta T+\frac1{u_j}.
\]
By the union bound, if $2k\delta<1$, then for all sufficiently large $T$,
\[
\left|[0,T]\setminus\bigcup_j B_j\right|
\ge(1-2k\delta)T-\sum_j u_j^{-1}>0.
\]
Removing the single point zero leaves a positive time with $F(t)\ge\delta$. Hence every real speed system has separation arbitrarily close to $1/(2k)$ from below. This statement deliberately does not assert attainment at the endpoint for an arbitrary real tuple: maximizing times can escape to infinity.

The conjecture asks for $1/(k+1)$, which is strictly larger than $1/(2k)$ when $k>1$. The loss is explicit: the union bound discards overlaps among the bad-time sets. Any improvement through this method must quantitatively exploit these overlaps for all possible relative speeds.

\paragraph{Endpoint attainment for rationally commensurable speeds.}
Suppose $u_j=cv_j$ for some $c>0$ and positive integers $v_j$. Rescale time by $c$. The resulting $F$ is continuous and one-periodic, so it attains a maximum on $[0,1]$. For each $\delta<1/(2k)$ the previous argument gives a point in this interval with $F\ge\delta$. Therefore its maximum is at least $1/(2k)$, now with attainment. The maximizing point is neither zero nor one because $F$ vanishes there. The case $k=1$ is even simpler: $t=1/(2u_1)$ gives the conjectured value $1/2$ for every real speed.

Dividing integer speeds by their common divisor does not change the maximum, since it only reparametrizes the circle variable. This reduction avoids treating large common factors as mathematically different cases. It does not bound the largest speed in a primitive tuple.

\paragraph{An exact congruence class at the desired threshold.}
If none of the integers $v_j$ is divisible by $k+1$, take $t=1/(k+1)$. The residue of each $v_j$ lies in $\{1,\ldots,k\}$ and consequently
\[
\left\|\frac{v_j}{k+1}\right\|\ge\frac1{k+1}.
\]
Thus this entire class satisfies the target. For example the speeds $1,\ldots,k$ attain the bound at that time.

For these particular speeds the bound is optimal. Among the $k+1$ points $0,t,2t,\ldots,kt$ on the circle, two cyclically consecutive points are within $1/(k+1)$, since their $k+1$ cyclic gaps sum to one. Their difference modulo one is $jt$ or $-jt$ for some $1\le j\le k$. Hence $F(t)\le1/(k+1)$ at every time. Repeated points cause a zero gap and only strengthen the argument. This proves sharpness of the proposed universal constant without proving universality.

\paragraph{Finite exact optimization for a fixed integer tuple.}
For $v_j\in\mathbb N$, the function $\|v_jt\|$ is piecewise affine on $[0,1]$, with breakpoints at $t=m/(2v_j)$. Partition the interval by the union of all these breakpoints. On each resulting interval every coordinate function is affine. Their minimum can change its active affine branch only at pairwise intersections. Consequently a maximum occurs at a partition endpoint or an intersection of two coordinate affine branches inside one partition interval (a constant active branch may be evaluated at an endpoint).

All candidate times and values are rational and can be compared exactly by integer arithmetic. Evaluating them gives a certificate for that tuple; rounding a decimal value near $1/(k+1)$ is unnecessary. A uniform proof would still have to cover infinitely many primitive integer tuples, and the number of intervals in this elementary algorithm depends on the magnitudes of the speeds.

\paragraph{Adversarial verification and remaining gap.}
The congruence argument fails when a speed is a multiple of $k+1$; this failure is of the chosen time, not a counterexample to the conjecture. Pairwise irrationality does not imply independent uniform phases for all runners, so no probabilistic independence assumption was made. Compactness of one period was used only for integer speeds, where a common period actually exists. Approximating a real tuple by rational ones would require control of the useful times to pass a closed inequality to a limit; this notebook does not assert such control.

No new finite computation is claimed. The exact first gap is to prove that the maximum separation is at least $1/(k+1)$ outside the congruence class, with uniform treatment of dependent real speeds. The measure argument gives a checked baseline and the piecewise-linear reduction specifies a rigorous way to test proposed integer counterexamples, but neither closes that gap.

\subsection{Sources}

{\small

\begin{itemize}

\item[\hypertarget{source:30:1}{S1}] ULAM.AI, \emph{UnsolvedMath}, supplied \texttt{problems.json}, record 30 (COMB-004), statement and background. The embedded literature-review date is 2026-08-17; that is the archive’s review date, not a fresh certification. Dataset landing page: \url{https://huggingface.co/datasets/ulamai/UnsolvedMath}.

\item[\hypertarget{source:30:2}{S2}] T. Sungkawichai and T. Trakulthongchai, Eleven, twelve, and thirteen lonely runners, arXiv:2604.23906 (2026). \url{https://arxiv.org/abs/2604.23906}. Bibliographic information imported from the supplied record.

\end{itemize}

}

\label{end:30}
\end{document}
